\section{Use exponent bits to select multiplications}
Write an exponent in binary. Each one bit selects a power of two: \(17=16+1\).
Rather than multiply the base seventeen times, square it to obtain powers one, two,
four, eight and sixteen, reducing each result modulo \(N\). Multiply the selected
powers into an accumulated result.

A right-to-left loop keeps three values: result \(r\), current base \(b\) and remaining
exponent \(e\). Initialize \(r=1\bmod N\), \(b=\text{base}\bmod N\) and \(e=E\).
If \(e\)'s low bit is one, replace \(r\) with \(rb\bmod N\). Shift \(e\) right by one.
If more bits remain, replace \(b\) with \(b^2\bmod N\) and repeat. The loop preserves
\(r b^e\equiv\text{base}^{E}\pmod N\). At \(e=0\), this leaves the desired result in
\(r\).

For base 65, exponent 17 and modulus 3233, all numbers in the table are decimal.
The last row skips an unnecessary square after the exponent becomes zero.
\tagpdfsetup{table/header-rows={1}}
\begin{center}
  \begin{tabular}{rrrrrr}
    \(r\) before & \(b\) before & \(e\) before & \(r\) after & \(b\) after & \(e\) after \\
    \hline
    1            & 65           & 17           & 65          & 992         & 8           \\
    65           & 992          & 8            & 65          & 1232        & 4           \\
    65           & 1232         & 4            & 65          & 1547        & 2           \\
    65           & 1547         & 2            & 65          & 789         & 1           \\
    65           & 789          & 1            & 2790        & 789         & 0           \\
  \end{tabular}
\end{center}

With \texttt{t0} holding \(b\), \texttt{t1} holding \(r\) and \texttt{t2} holding \(e\),
the conditional multiply has a short observation boundary. The modulus remains in
\texttt{a2}.

\Needspace{11\baselineskip}
\MARSCodeLabel{Select the current base using the remaining exponent	extquotesingle s low bit.}
\begin{minted}[samepage=true]{mars-idea}
andi $t3, $t2, 1
    beqz $t3, pow32_shift
    nop
    dmulu $t1, $t1, $t0
    dmodu $t1, $t1, $a2
\end{minted}

After the remainder, \texttt{t1} is again a residue below \(N\). Inspect it at this
point, before a later call or instruction reuses a temporary register. An exponent
of zero must reach the output without entering this multiply; initialize the identity
before testing the loop condition.
